\subsection{Derivatives of Basic Functions}

\begin{prop}
  Let $f : D \rightarrow \mathbb{R}, f(x) = x^2$ where $D = \mathbb{R}$ if $n > 0$ and $\mathbb{R} \setminus \{0\}$ otherwise. Then $f$ is differentiable on $D$ with derivative
  \[f'(x) = n x ^{n-1}.\]
\end{prop}
\begin{proof}
    \begin{itemize}
      \item Suppose $n \geq 0$. Then using induction on $n$, which is true $n = 0$, then assuming true for $n= k$, for $n = k + 1$
        \begin{align*}
          f'(x) &= \frac{d}{dx} (x)(x)^k  \\
                &= x^k + (k)x^{k-1} x \\
                &= (k + 1)x^k
        .\end{align*}
      \item Suppose $n < 0$. Then $f = \frac{1}{g}$ where $g(x) = x^{-n}$. Then using the quotient rule
        \begin{align*}
                &= - \frac{g'(x)}{g(x)^2} \\
                &= n x^{n - 1}
        .\end{align*}
    \end{itemize}
\end{proof}

\begin{corollary}
Let $p: \mathbb{R} \rightarrow \mathbb{R}$ be a polynomial with 
\[P(x) = a_n x^n + ... + a_1 x + a_0.\]
Then $P$ is differentiable on $\mathbb{R}$ and
\[p'(x) = n a_n x^{n - 1} + ... + a_1.\]
\end{corollary}

\begin{prop}
    The functions $\exp, \cos, \sin$ are differentiable with
    \begin{itemize}
      \item $\exp'(x) = \exp(x)$
      \item $\cos'(x) = - \sin (x)$
      \item $\sin'(x) = \cos(x)$
    \end{itemize}
\end{prop}
\begin{proof}
   The arguments for all three are similar, so let the desired function be $f : \mathbb{R} \rightarrow \mathbb{R}$.
   \begin{itemize}
     \item Step 1: Show $f$ is differentiable at $0$ and find $f'(0)$.
     \item Step 2: In the derivative of $f(x)$, write $f(x + h)$ as a product of $f(h)$, then as this is already known it can be evaluated.
   \end{itemize}
\end{proof}

\subsection{The Chain Rule}
\begin{theorem}
    Let $D, E \subseteq \mathbb{R}$, $f : D \rightarrow \mathbb{R}$ and $g : E \rightarrow D$. Let $c \in E$ be an interior point, such that $g(c)$ is an interior point of $D$. If
    \begin{itemize}
      \item $g$ is differentiable at $c$
      \item $f$ is differentiable at $g(c)$      \item 
    \end{itemize}
     then $f \circ g$ is differentiable at $c$ with
     \[(f \circ g)'(c) = f'(g(c))g'(c).\]

\end{theorem}
\begin{proof}
  Let $h \in \mathbb{R} \setminus \{0\}$ be such that $c + h \in E$. If $g(c + h) \not = g(c)$ then
  \begin{align*}
    \frac{(f \circ g)(c + h) - (f \circ g)(c)}{h} &= \frac{f(g(c + h)) - f(g(c))}{h} \\
      &= \frac{f(g(c+h)) - f(g(c))}{g(c + h) - g(c)} \cdot \frac{g (c+h)- g(c)}{h}
  .\end{align*}
  We need to define $q : D \rightarrow \mathbb{R}$ by
  \[q(y) = \left\{\begin{array}{lr}
      \frac{f(y) - f(g(c))}{y - g(c)} & y \in D \setminus \{g(c)\}\\
  f'(g(c)) & y = g(c)\\
  \end{array} \right.\]
  which lets us more easily prove properties needed. Then, 
  \begin{align*}
    \frac{(f \circ g)(c + h) - (f \circ g)(c)}{h} &= q(g(c + h)) \cdot \frac{g(c + h) - g(c)}{h}
  \end{align*}
  which is true for both cases of $q$, letting us avoid the case $g(c+h) = g(c)$. Then 
  \[\lim_{y \rightarrow g(c)} q(y) = f'(g(c)) = q(g(c))\]
  so $q$ is continuous, and so is $q \circ g$, therefore
  \begin{align*}
    ( f \circ g)'(c) &= \lim_{h \rightarrow 0} q(g(c + h)) \lim_{h \rightarrow 0} \frac{g(c + h) - g(c)}{h} \\
                     &= \lim_{h \rightarrow 0} q(g(c + h)) g'(c) \\
                     &= f'(g(c)) g'(c)
  .\end{align*}
\end{proof}

\subsection{The derivative of inverse functions}


\begin{theorem}
  Let $D, E \subseteq \mathbb{R}$ and let $f : D \rightarrow E$ bijective function. Suppose $c \in D$ and $f(c) \in E$ are interior points. If $f$ is differentiable at $c$ with $f'(c) \not = 0$ then $f^{-1}$ is differentiable at $f(c)$ with
  \[(f^{-1})'(f(c)) = \frac{1}{f'(c)}.\]
\end{theorem}
\begin{proof}
    Non examinable.
\end{proof}

